% ========================================================================================
% CAPITOLO 7 - SWITCH MULTISTADIO
% ========================================================================================
\chapter{Switch multistadio}

% ----------------------------------------------------------------------------------------
\section{Esercizio 11: Switch multistadio}
% ----------------------------------------------------------------------------------------

\subsection{Esercizio 11.1}

\begin{traccia}
	Progettare ed implementare in VHDL uno switch multistadio secondo il modello omega network.
	Lo switch deve consentire lo scambio di messaggi di 2 bit ciascuno da un nodo sorgente a un
	nodo destinazione in una rete con 4 nodi, implementando uno schema a priorità fissa fra i
	nodi (es. nodo 1 più prioritario, con priorità decrescenti fino al nodo 4).
\end{traccia}

\progetto

La rete Omega con $N=4$ nodi ha $\log_2 N = 2$ stadi di $N/2 = 2$ switch elementari
$2\times2$, preceduti da una permutazione \emph{perfect shuffle}
(Figura~\ref{fig:omega4}). Allo stadio $k$ ogni switch instrada il messaggio in base al
$k$-esimo bit dell'indirizzo di destinazione, a partire dal più significativo:
\texttt{0} verso l'uscita superiore, \texttt{1} verso quella inferiore.

\begin{figure}[H]
	\centering
	\begin{tikzpicture}[x=1cm, y=0.8cm]
		% switch
		\foreach \s/\x in {0/2.5, 1/6} {
			\foreach \k in {0,1} {
				\node[blocco, minimum width=1.1cm, minimum height=1.6cm] (sw\s\k) at (\x,{-2.5*\k-0.75}) {\scriptsize SW\,\s,\k};
			}
		}
		% ingressi: sorgente -> shuffle -> stadio 0 (0,2 -> SW00; 1,3 -> SW01)
		\foreach \src/\y in {0/0, 1/-1, 2/-2.5, 3/-3.5} {
			\node[left, font=\small] (in\src) at (0,\y) {sorg.\,\src};
		}
		\draw[seg] (0,0)    -- (1.2,0)    -- (1.6,-0.4) -- (sw00.west |- 0,-0.4);
		\draw[seg] (0,-2.5) -- (1.2,-2.5) -- (1.6,-1.1) -- (sw00.west |- 0,-1.1);
		\draw[seg] (0,-1)   -- (1.2,-1)   -- (1.6,-2.9) -- (sw01.west |- 0,-2.9);
		\draw[seg] (0,-3.5) -- (1.2,-3.5) -- (1.6,-3.6) -- (sw01.west |- 0,-3.6);
		% shuffle fra gli stadi: sup -> SW10, inf -> SW11
		\draw[seg] (sw00.east |- 0,-0.4) -- (sw10.west |- 0,-0.4);
		\draw[seg] (sw00.east |- 0,-1.1) -- (4.1,-1.1) -- (4.6,-2.9) -- (sw11.west |- 0,-2.9);
		\draw[seg] (sw01.east |- 0,-2.9) -- (4.1,-2.9) -- (4.6,-1.1) -- (sw10.west |- 0,-1.1);
		\draw[seg] (sw01.east |- 0,-3.6) -- (sw11.west |- 0,-3.6);
		% uscite
		\foreach \d/\y in {0/-0.4, 1/-1.1, 2/-2.9, 3/-3.6} {
			\draw[seg] (sw10.east |- 0,\y) -- (8,\y) node[right, font=\small]{dest.\,\d};
		}
		\node[font=\sffamily\scriptsize, black!60] at (2.5,1) {stadio 0 (bit $d_1$)};
		\node[font=\sffamily\scriptsize, black!60] at (6,1) {stadio 1 (bit $d_0$)};
	\end{tikzpicture}
	\caption{Rete Omega a 4 nodi: due stadi di switch $2\times 2$ collegati da perfect shuffle.}
	\label{fig:omega4}
\end{figure}

\begin{osservazione}[Rete bloccante]
	La rete Omega è bloccante: due messaggi con destinazioni diverse possono contendersi lo
	stesso collegamento interno. Ad esempio, sorgente 0 $\rightarrow$ destinazione 0 e
	sorgente 2 $\rightarrow$ destinazione 1 richiedono entrambe l'uscita superiore di
	SW\,0,0. La priorità fissa stabilisce quale dei due messaggi passa.
\end{osservazione}

\todo{Interfaccia dei nodi, arbitro a priorità fissa, una o più trasmissioni per volta.}

\implementazione
\todo{Codice di \ent{switch2x2}, dell'arbitro e della rete.}

\simulazione
\todo{Messaggi verso ogni destinazione, messaggi in conflitto.}

\subsection{Esercizio 11.2 (solo 9 CFU)}

\begin{traccia}
	Rimuovendo l'ipotesi di lavorare secondo uno schema a priorità fissa fra i nodi e
	considerando una rete di 8 nodi, lo switch deve gestire eventuali conflitti generati da
	collisioni con un meccanismo a scelta dello studente.
\end{traccia}

\todo{Rete a 3 stadi, meccanismo di risoluzione dei conflitti.}

\subsection{Esercizio 11.3 (solo 9 CFU)}

\begin{traccia}
	Si implementi un protocollo di handshaking semplice regolato da una coppia di segnali
	(pronto a inviare/pronto a ricevere) per l'invio di ciascun messaggio fra due nodi.
\end{traccia}

\todo{Protocollo e sua integrazione nella rete.}
